\scanpage{48}
\mspage{135}

à source non vide, et à but connexe, est surjective
(\uline{NB} \unsure{toutes} les fibres d'une $=$ comp. connexe
de $S$ sont Hot-équivalentes\deleted{\unsure{)}}, donc si l'une est non
vide, toutes le sont.) Cela achève la démonstration, modulo le lemme.

Il reste à prouver celui-ci. Comme $S=\varnothing$, je
peux choisir une ``origine''
\[
s_0\in S,
\]
et regarder la catégorie
\[
\underline{\mathrm{Ch}}_{\infty}(\deleted{\unsure{s_0|s}}S;s_0,-)
\]
des chemins infinis dans $S$, d'origine $s_0$, fibre
de $\unsure{\underline{\mathrm{Ch}}|\underline{\mathcal A}}_{\infty}(S)$ par le foncteur-source \deleted{\unsure{$\sigma$}}
\[
\sigma:\underline{\mathrm{Ch}}_{\infty}(S)\longrightarrow S.
\]
On sait que les fibres de $\sigma$ sont Hot-asphériques,
(VII), et la catégorie construite ici n'est
autre \deleted{\unsure{que}} (avec les notations de \uline{$\overline{\mathrm{VII}}$} p.152) que
la \deleted{\unsure{fo}} cat. $\overline{S'}$ qui entre dans la factorisation
canonique de
\[
f:\{s_0\}=S'\unsure{\longrightarrow|\text{---}}S
\]
en
\[
\{s_0\}=S'\xrightarrow[\text{Hot-équiv.}]{}\overline{S'}
 \xrightarrow{\bar f} S,
\]
où ici $\bar f$ s'interprète simplement comme
le foncteur but (i.e le foncteur induit sur

% archive sheet number: 48
